\documentclass[11pt,a4paper]{article}

\usepackage{fontspec}
\usepackage[margin=22mm]{geometry}
\usepackage{amsmath,amssymb}
\usepackage{booktabs,longtable,array,tabularx}
\usepackage{graphicx}
\usepackage{caption}
\usepackage{xcolor}
\usepackage{hyperref}
\usepackage{enumitem}

\defaultfontfeatures{Ligatures=TeX,Scale=MatchLowercase}
\IfFontExistsTF{Noto Sans CJK KR}{
  \setmainfont{Noto Sans CJK KR}[
    ItalicFeatures={FakeSlant=0.18},
    BoldItalicFeatures={FakeSlant=0.18}
  ]
  \setsansfont{Noto Sans CJK KR}[
    ItalicFeatures={FakeSlant=0.18},
    BoldItalicFeatures={FakeSlant=0.18}
  ]
}{
  \IfFontExistsTF{Apple SD Gothic Neo}{
    \setmainfont{Apple SD Gothic Neo}[
      ItalicFeatures={FakeSlant=0.18},
      BoldItalicFeatures={FakeSlant=0.18}
    ]
    \setsansfont{Apple SD Gothic Neo}[
      ItalicFeatures={FakeSlant=0.18},
      BoldItalicFeatures={FakeSlant=0.18}
    ]
  }{
    \setmainfont{NanumGothic}[
      ItalicFeatures={FakeSlant=0.18},
      BoldItalicFeatures={FakeSlant=0.18}
    ]
    \setsansfont{NanumGothic}[
      ItalicFeatures={FakeSlant=0.18},
      BoldItalicFeatures={FakeSlant=0.18}
    ]
  }
}
\IfFontExistsTF{DejaVu Sans Mono}{
  \setmonofont{DejaVu Sans Mono}
}{
  \IfFontExistsTF{Menlo}{
    \setmonofont{Menlo}
  }{
    \setmonofont{Liberation Mono}
  }
}
\XeTeXlinebreaklocale "ko"
\XeTeXlinebreakskip = 0pt plus 1pt

\hypersetup{
  colorlinks=true,
  linkcolor=blue!50!black,
  urlcolor=blue!50!black,
  citecolor=blue!50!black,
  pdftitle={MuJoCo Fluid Force Mapping v3.0},
  pdfauthor={OpenAI Codex}
}

\setlength{\parskip}{0.55em}
\setlength{\parindent}{0pt}
\setlength{\emergencystretch}{3em}
\renewcommand{\arraystretch}{1.22}
\newcolumntype{Y}{>{\raggedright\arraybackslash}X}
\newcommand{\code}[1]{\nolinkurl{#1}}
\captionsetup{
  font=small,
  labelfont=bf,
  justification=raggedright,
  singlelinecheck=false,
  skip=6pt
}

\title{MuJoCo Fluid Force 공식식과 v3.0 매핑 정리\\
\large MuJoCo 공식 문서 기준으로 각 항의 물리적 의미와 현재 UUV v3.0 구현 대응관계를 정리한 별도 문서}
\author{조백규 교수님 연구실 최애의 제자 박사과정 보고용 정리}
\date{2026-04-06}

\begin{document}
\maketitle
\tableofcontents
\clearpage

\section{문서 목적과 기준}

이 문서는 \textbf{MuJoCo 공식 문서의 fluid force 정의}를 기준으로, 현재 \code{uuv\_mujoco/v2.2} 의 \textbf{v3.0 active current model} 이 그 공식을 어디까지 그대로 쓰고, 어디부터는 runtime에서 별도로 보정하는지를 정리한 별도 보고서다.

이 문서가 답하려는 질문은 네 가지다.

\begin{enumerate}[leftmargin=2em]
  \item MuJoCo의 \textbf{inertia-based fluid model} 과 \textbf{ellipsoid-based fluid model} 은 각각 어떤 물리식으로 구성되는가.
  \item ellipsoid 모델의 각 항 \(\mathbf{f}_A, \mathbf{f}_D, \mathbf{f}_M, \mathbf{f}_K, \mathbf{f}_V, \mathbf{g}_A, \mathbf{g}_D, \mathbf{g}_V\) 는 무슨 의미를 가지는가.
  \item XML의 \code{density}, \code{viscosity}, \code{fluidshape}, \code{fluidcoef} 가 실제로 어떤 항을 바꾸는가.
  \item 현재 v3.0 current runtime에서 MuJoCo built-in fluid term과 Python runtime 보정 항이 어디서 갈리는가.
\end{enumerate}

\textbf{공식 기준 문서}는 다음 두 페이지다.

\begin{itemize}[leftmargin=2em]
  \item MuJoCo official documentation, \href{https://mujoco.readthedocs.io/en/3.1.1/computation/fluid.html}{Fluid forces}
  \item MuJoCo official documentation, \href{https://mujoco.readthedocs.io/en/3.1.1/XMLreference.html}{XML Reference}
\end{itemize}

즉 이 문서의 공식식은 임의의 블로그나 2차 자료가 아니라 \textbf{MuJoCo 공식 문서에 나온 정의와 식}을 기준으로 다시 정리한 것이다. 여기서 \textbf{현재 repository 구현 설명}은 그 위에 얹은 코드 매핑 해설이다.

\section{MuJoCo fluid model의 큰 구조}

MuJoCo는 물속 유동장을 풀지 않는다. 대신 공식 문서가 명시하듯이 \textbf{state-less phenomenological model} 을 제공한다. 즉 주변 유체에 별도 상태 변수를 두지 않고, body 속도와 형상 근사만으로 수중 힘과 토크를 근사한다.

공식 문서 기준으로 fluid model은 두 종류다.

\begin{enumerate}[leftmargin=2em]
  \item \textbf{Inertia model}: body 관성과 질량에서 equivalent inertia box를 만들고, 그 박스를 기준으로 drag와 viscous resistance를 계산한다.
  \item \textbf{Ellipsoid model}: 각 geom을 ellipsoid로 근사하고, added mass, drag, Magnus lift, Kutta lift, viscous resistance까지 포함한 더 세분화된 fluid force를 계산한다.
\end{enumerate}

두 모델 모두 medium의 \textbf{density \(\rho\)} 와 \textbf{kinematic viscosity \(\beta\)} 를 사용한다. 다만 ellipsoid model은 여기에 더해 \textbf{geom별 \code{fluidcoef} 다섯 개}를 받는다.

\subsection{현재 v3.0 current path에서 active한 모델}

현재 active current path는 \code{tank\_current\_scene.xml} 에 있는 \code{fluid\_*} geom에 \code{fluidshape="ellipsoid"} 와 \code{fluidcoef} 를 정의하고, \code{run\_urdf\_full.py} 에서 \textbf{ellipsoid-only dynamics} 를 선택하는 구조다. 이 상태에서는:

\begin{itemize}[leftmargin=2em]
  \item MuJoCo built-in ellipsoid fluid가 active이다.
  \item 공식 문서가 말하는 \textbf{body inertia sizes 기반 inertia model은 parent body에서 비활성화}된다.
  \item Python runtime의 예전 coefficient-based 6-DOF custom hydrodynamics는 current path에서 끈다.
  \item 대신 \textbf{buoyancy, distributed mass/inertia synthesis, thruster shaping, reaction torque, sensor export} 는 여전히 Python runtime이 담당한다.
\end{itemize}

즉 \textbf{현재 모델은 ``순수 MuJoCo만''도 아니고 ``순수 Python만''도 아니다.} fluid term은 MuJoCo ellipsoid model을 주력으로 쓰고, hydrostatic/actuation/interface는 runtime이 보강하는 하이브리드 구조다.

\section{Inertia model 공식식}

공식 문서에 따르면 inertia model은 body 질량 \(\mathcal{M}\) 과 관성행렬 \(\mathcal{I}\) 로부터 equivalent inertia box 반치수 \((r_x, r_y, r_z)\) 를 만든다.

\begin{align*}
r_x &= \sqrt{\frac{3}{2 \mathcal{M}}\left(\mathcal{I}_{yy} + \mathcal{I}_{zz} - \mathcal{I}_{xx}\right)} \\
r_y &= \sqrt{\frac{3}{2 \mathcal{M}}\left(\mathcal{I}_{zz} + \mathcal{I}_{xx} - \mathcal{I}_{yy}\right)} \\
r_z &= \sqrt{\frac{3}{2 \mathcal{M}}\left(\mathcal{I}_{xx} + \mathcal{I}_{yy} - \mathcal{I}_{zz}\right)}
\end{align*}

그 다음 body frame 선속도 \(\mathbf{v}\) 와 각속도 \(\boldsymbol{\omega}\) 를 이용해

\begin{align*}
\mathbf{f}_{\text{inertia}} &= \mathbf{f}_D + \mathbf{f}_V \\
\mathbf{g}_{\text{inertia}} &= \mathbf{g}_D + \mathbf{g}_V
\end{align*}

를 계산한다. 여기서 \(D\) 는 quadratic drag, \(V\) 는 viscous resistance다.

\subsection{Quadratic drag}

\begin{align*}
f_{D,i} &= - 2 \rho r_j r_k |v_i| v_i \\
g_{D,i} &= - \frac{1}{2}\rho r_i \left(r_j^4 + r_k^4\right)|\omega_i|\omega_i
\end{align*}

의미는 단순하다.

\begin{itemize}[leftmargin=2em]
  \item 선속도 크기가 커질수록 힘이 \(|v|v\) 로 커진다.
  \item body가 어떤 축으로 넓게 보이느냐가 \(r_j r_k\) 항에 들어간다.
  \item 회전 저항은 회전축 주위 swept area 효과가 \(\left(r_j^4 + r_k^4\right)\) 꼴로 들어가서, 긴 body는 yaw/roll/pitch damping이 크게 달라질 수 있다.
\end{itemize}

\subsection{Viscous resistance}

\begin{align*}
f_{V,i} &= - 6 \beta \pi r_{eq} v_i \\
g_{V,i} &= - 8 \beta \pi r_{eq}^3 \omega_i
\end{align*}

여기서 \(r_{eq} = (r_x + r_y + r_z)/3\) 이다. 이 항은 \textbf{저 Reynolds 영역의 선형 저항}에 해당한다.

\textbf{영향}:
\begin{itemize}[leftmargin=2em]
  \item density와 달리 viscosity만 올려도 선형 damping이 증가한다.
  \item 속도가 아주 작을 때 잔류 진동이나 꼬리 흔들림을 정리하는 데 효과가 크다.
  \item 현재 프로젝트에서 ``조종감이 너무 미끄럽다''를 잡을 때 viscosity나 선형 damping 계열이 중요한 이유가 이 항 때문이다.
\end{itemize}

\section{Ellipsoid model 공식식}

\subsection{총합 구조}

공식 문서 기준 ellipsoid model은 다음처럼 분해된다.

\begin{align*}
\mathbf{f}_{\text{ellipsoid}} &= \mathbf{f}_A + \mathbf{f}_D + \mathbf{f}_M + \mathbf{f}_K + \mathbf{f}_V \\
\mathbf{g}_{\text{ellipsoid}} &= \mathbf{g}_A + \mathbf{g}_D + \mathbf{g}_V
\end{align*}

\begin{center}
\small
\begin{tabularx}{\linewidth}{>{\raggedright\arraybackslash}p{0.16\linewidth}YY}
\toprule
항 & 물리적 의미 & 현재 tuning에서 체감되는 효과 \\
\midrule
\(\mathbf{f}_A, \mathbf{g}_A\) & added mass / added inertia & 급가속, 급회전에서 ``물까지 같이 끌고 가는'' 무거운 느낌 \\
\(\mathbf{f}_D, \mathbf{g}_D\) & quadratic drag / angular drag & 속도가 커질수록 강하게 붙는 감쇠, 회전 꼬리 억제 \\
\(\mathbf{f}_M\) & Magnus lift & 회전하면서 전진할 때 옆으로 휘는 coupling \\
\(\mathbf{f}_K\) & Kutta lift & 받음각이 생길 때 양력성 측력 발생 \\
\(\mathbf{f}_V, \mathbf{g}_V\) & linear viscous resistance & 저속 잔류 motion 정리, 점성 damping \\
\bottomrule
\end{tabularx}
\end{center}

\subsection{Projected area}

ellipsoid drag와 lift 식에는 projected area가 반복해서 들어간다. 공식 문서가 제시하는 projected area는

\[
A^{\mathrm{proj}}_{\mathbf{u}}
=
\pi
\sqrt{
\frac{
r_y^4 r_z^4 u_x^2 + r_z^4 r_x^4 u_y^2 + r_x^4 r_y^4 u_z^2
}{
r_y^2 r_z^2 u_x^2 + r_z^2 r_x^2 u_y^2 + r_x^2 r_y^2 u_z^2
}}
\]

이다. 여기서 \(\mathbf{u}\) 는 속도 방향 단위벡터다.

\textbf{의미}:
\begin{itemize}[leftmargin=2em]
  \item 속도 방향이 바뀌면 projected area도 바뀐다.
  \item 같은 ellipsoid라도 broadside로 움직일 때와 slender 방향으로 움직일 때 drag가 달라진다.
  \item 결국 \textbf{geom size와 geom orientation} 이 fluid force 방향성과 크기를 동시에 결정한다.
\end{itemize}

\subsection{Added mass}

공식 문서는 added mass를 potential flow 관점으로 설명하고, 세 대칭면을 가진 body에 대해 다음 compact form을 제시한다.

\begin{align*}
\mathbf{f}_A &= - \mathbf{m}_A \circ \dot{\mathbf{v}} + (\mathbf{m}_A \circ \mathbf{v}) \times \boldsymbol{\omega} \\
\mathbf{g}_A &= - \mathbf{I}_A \circ \dot{\boldsymbol{\omega}} + (\mathbf{m}_A \circ \mathbf{v}) \times \mathbf{v} + (\mathbf{I}_A \circ \boldsymbol{\omega}) \times \boldsymbol{\omega}
\end{align*}

즉 added mass는 단순 damping이 아니라 \textbf{가속도와 속도-회전 coupling} 에 반응하는 항이다.

ellipsoid의 added mass 계수는

\begin{align*}
\kappa_i &= \int_0^\infty
\frac{r_i r_j r_k}
{\sqrt{(r_i^2 + \lambda)^3 (r_j^2 + \lambda)(r_k^2 + \lambda)}}
\mathrm{d}\lambda \\
m_{A,i} &= \rho V \frac{\kappa_i}{2 - \kappa_i}
\end{align*}

로 주어진다. \(I_{A,i}\) 도 공식 문서에 별도 식으로 정의된다.

\textbf{영향}:
\begin{itemize}[leftmargin=2em]
  \item constant velocity에는 직접 힘을 남기지 않지만, \textbf{속도 변화에 저항}한다.
  \item body가 갑자기 방향을 틀 때 ``즉시 따라가지 않고 묵직하게 버티는'' 느낌은 added mass 성분이 핵심이다.
  \item current path에서 이 항은 MuJoCo ellipsoid model이 맡고, Python custom added-mass term은 끈다.
\end{itemize}

\subsection{Quadratic translational drag}

공식 문서는 translational drag를

\[
\mathbf{f}_D =
- \rho
\left[
C_{D,\text{blunt}} A^{\mathrm{proj}}_{\mathbf{v}}
+
C_{D,\text{slender}}\left(A_{\max} - A^{\mathrm{proj}}_{\mathbf{v}}\right)
\right]
\|\mathbf{v}\| \mathbf{v}
\]

로 쓴다. 여기서 \(A_{\max} = 4 \pi r_{\max} r_{\min}\) 이다.

\textbf{중요한 해석}:
\begin{itemize}[leftmargin=2em]
  \item \(C_{D,\text{blunt}}\): body를 정면으로 받았을 때의 broadside drag를 키운다.
  \item \(C_{D,\text{slender}}\): 길쭉한 방향으로 흐를 때의 drag baseline을 키운다.
  \item 두 항을 섞기 때문에, 하나의 계수만으로는 안 되는 ``자세에 따른 저항 차이'' 를 표현할 수 있다.
\end{itemize}

\subsection{Angular drag}

공식 문서는 angular drag reference inertia를

\[
\mathbf{I}_{D,ii} = \frac{8\pi}{15} r_i \max(r_j, r_k)^4
\]

로 두고, torque를

\[
\mathbf{g}_D
=
- \rho \boldsymbol{\omega}
\Big(
[C_{D,\text{angular}}\mathbf{I}_D + C_{D,\text{slender}}(\mathbf{I}_{\max}-\mathbf{I}_D)]
\cdot
\boldsymbol{\omega}
\Big)
\]

로 계산한다.

\textbf{영향}:
\begin{itemize}[leftmargin=2em]
  \item \(C_{D,\text{angular}}\) 는 roll/pitch/yaw 자체의 감쇠를 직접 키운다.
  \item 회전 후 계속 더 도는 tail motion을 줄이는 데 가장 직접적인 knob다.
  \item current scene에서 yaw가 지나치게 남는다면 fluidcoef의 세 번째 값과 geom size가 우선 검토 대상이다.
\end{itemize}

\subsection{Linear viscous resistance}

\begin{align*}
\mathbf{f}_V &= -6\pi r_D \beta \mathbf{v} \\
\mathbf{g}_V &= -8\pi r_D^3 \beta \boldsymbol{\omega}
\end{align*}

여기서 \(r_D = (r_x+r_y+r_z)/3\) 이다.

\textbf{영향}:
\begin{itemize}[leftmargin=2em]
  \item 저속 damping, 미세 진동, 잔류 속도 정리.
  \item viscosity \(\beta\) 가 커질수록 효과가 커진다.
  \item density를 안 바꾸고도 ``조금 더 점성 있게'' 만들 수 있다.
\end{itemize}

\subsection{Magnus lift}

공식 문서는 Magnus force를

\[
\mathbf{f}_M = C_M \rho V (\boldsymbol{\omega}\times\mathbf{v})
\]

로 근사한다.

\textbf{영향}:
\begin{itemize}[leftmargin=2em]
  \item 회전하면서 전진할 때 lateral force가 생긴다.
  \item yaw rate와 surge/sway가 동시에 있을 때 \textbf{옆으로 밀리거나 기묘한 arc path} 를 만들 수 있다.
  \item 수중 ROV에서는 항공 곤충만큼 지배적이지 않을 수 있지만, spin-translation coupling을 설명하는 항이다.
\end{itemize}

\subsection{Kutta lift}

공식 문서는 Kutta condition 기반 lift를

\begin{align*}
\mathbf{f}_K
&=
\frac{C_K \rho A^{\mathrm{proj}}_{\mathbf{v}}}{\|\mathbf{v}\|}
\big(\mathbf{v}\times\mathbf{v}_{\parallel}\big)\times\mathbf{v} \\
&=
C_K \rho A^{\mathrm{proj}}_{\mathbf{v}}
\left(\hat{\mathbf{v}}\cdot\hat{\mathbf{n}}_{s,\mathbf{v}}\right)
\left(\hat{\mathbf{n}}_{s,\mathbf{v}} \times \mathbf{v}\right)\times\mathbf{v}
\end{align*}

로 제시한다.

\textbf{영향}:
\begin{itemize}[leftmargin=2em]
  \item 받음각이나 cross-flow가 있을 때 lift-like side force를 만든다.
  \item body가 완전 구형이면 이 항은 0으로 간다.
  \item 판형/평판형 body, 혹은 특정 방향으로 납작한 body에서 더 의미가 커진다.
\end{itemize}

\section{XML level에서 무엇을 바꾸면 어떤 항이 바뀌는가}

\begin{center}
\small
\begin{tabularx}{\linewidth}{>{\raggedright\arraybackslash}p{0.18\linewidth}YY}
\toprule
XML knob & 공식 모델에서 바뀌는 것 & 현재 v3.0에서의 실질적 의미 \\
\midrule
\code{option density} & \(\rho\) 전체 스케일 & drag, added mass, lift 전반 스케일 변화 \\
\code{option viscosity} & \(\beta\) 전체 스케일 & 저속 선형 damping, 회전 점성 저항 증가 \\
\code{geom fluidshape="ellipsoid"} & ellipsoid model 활성화 & 해당 body에서 inertia model 비활성화, geom-level fluid 활성화 \\
\code{geom fluidcoef[0]} & \(C_{D,\text{blunt}}\) & broadside drag 증가 \\
\code{geom fluidcoef[1]} & \(C_{D,\text{slender}}\) & slender-axis drag와 angular drag 혼합 증가 \\
\code{geom fluidcoef[2]} & \(C_{D,\text{angular}}\) & yaw/roll/pitch 회전 감쇠 증가 \\
\code{geom fluidcoef[3]} & \(C_K\) & cross-flow lift 증가 \\
\code{geom fluidcoef[4]} & \(C_M\) & spin-translation Magnus coupling 증가 \\
\code{geom size} & \(r_x,r_y,r_z\) 및 projected area & 어떤 축에서 저항과 added mass가 큰지 결정 \\
\code{geom pos} & fluid center 위치 & 같은 drag여도 CoM 기준 lever arm이 달라져 torque 체감이 달라짐 \\
\bottomrule
\end{tabularx}
\end{center}

\section{현재 v3.0 current scene의 실제 값}

현재 active scene은 \code{uuv\_mujoco/v2.2/scenes/tank\_current\_scene.xml} 이다. 여기서 fluid proxy geom은 세 개다.

\begin{center}
\small
\begin{tabularx}{\linewidth}{>{\raggedright\arraybackslash}p{0.28\linewidth}Y>{\raggedright\arraybackslash}p{0.20\linewidth}>{\raggedright\arraybackslash}p{0.25\linewidth}}
\toprule
geom name & position [m] & size [m] & fluidcoef \\
\midrule
\code{fluid\_center\_enclosure} & \shortstack[l]{(-0.005, 0.000,\\ 0.020)} & \((0.205,\ 0.115,\ 0.130)\) & \shortstack[l]{(0.620,\ 0.175,\\ 1.800,\ 0.0,\ 0.0)} \\
\code{fluid\_port\_lower\_body} & \shortstack[l]{(0.000,\ 0.208,\\ 0.020)} & \((0.370,\ 0.080,\ 0.120)\) & \shortstack[l]{(0.520,\ 0.150,\\ 1.300,\ 0.0,\ 0.0)} \\
\code{fluid\_starboard\_lower\_body} & \shortstack[l]{(0.000,\ -0.208,\\ 0.020)} & \((0.370,\ 0.080,\ 0.120)\) & \shortstack[l]{(0.520,\ 0.150,\\ 1.300,\ 0.0,\ 0.0)} \\
\bottomrule
\end{tabularx}
\end{center}

이 값을 읽으면 중요한 사실이 세 가지 보인다.

\begin{enumerate}[leftmargin=2em]
  \item \textbf{center enclosure} 가 angular drag 계수 \(1.8\) 로 가장 크다. 즉 중앙 body가 회전 감쇠를 더 강하게 준다.
  \item 좌우 lower body 두 개는 같은 계수를 가지므로, 좌우 비대칭 yaw torque를 만들기보다 \textbf{대칭적인 lateral / yaw damping scaffold} 역할을 한다.
  \item \(C_K\) 와 \(C_M\) 는 현재 scene에서 \textbf{0} 이다. 즉 current v3.0 active scene은 Kutta lift와 Magnus lift를 적극적으로 쓰는 설계가 아니라, \textbf{added mass + drag + angular drag + viscous resistance 중심} 의 conservative 설정이다.
\end{enumerate}

\section{current profile에서 MuJoCo fluid와 함께 바뀌는 값}

MuJoCo fluid model만으로 현재 응답을 다 설명할 수는 없다. current profile에는 다음 값들이 함께 들어 있다.

\begin{center}
\small
\begin{tabularx}{\linewidth}{>{\raggedright\arraybackslash}p{0.26\linewidth}>{\raggedright\arraybackslash}p{0.14\linewidth}>{\raggedright\arraybackslash}p{0.14\linewidth}Y}
\toprule
항목 & legacy & current & 현재 체감에 주는 의미 \\
\midrule
\code{buoyancy\_scale} & 1.003 & 1.008 & current가 약간 더 강한 buoyancy margin을 가진다. \\
\code{surface\_heave\_damping} & 18.0 & 0.0 & legacy의 수면 근처 추가 heave damping이 current에서는 제거됐다. \\
\code{cob\_torque\_scale} & 0.2 & 1.0 & current는 CoB restoring torque를 더 직접적으로 살린다. \\
\code{cob\_x\_offset} & -0.003 & 0.009 & fluid/부력 기준점이 legacy와 current에서 다르다. \\
\code{thruster\_voltage} & 20.0 & 16.0 & 추력 authority baseline이 current에서 더 보수적이다. \\
\code{thruster\_force\_max} & 45.0 & 21.0 & current는 thrust ceiling 자체가 더 낮다. \\
\code{linear\_drag} & 1.10 & 1.55 & current는 저항 baseline이 더 크다. \\
\code{angular\_drag} & 0.32 & 0.48 & current는 회전 감쇠 baseline도 더 크다. \\
\code{body\_inertia\_scale\_xyz} & (1,1,1) & [1e-6, 1e-6, 1e-6] & current는 body component 합성 관성을 사실상 극단적으로 줄인다. \\
\bottomrule
\end{tabularx}
\end{center}

\textbf{여기서 핵심은} current 응답 안정화가 MuJoCo fluidcoef 하나만으로 생긴 것이 아니라는 점이다. 실제 current v3.0의 안정적이고 보수적인 거동은

\begin{enumerate}[leftmargin=2em]
  \item built-in ellipsoid fluid 활성화,
  \item custom 6-DOF hydrodynamics 비활성화,
  \item 큰 drag / angular drag baseline,
  \item 매우 작은 runtime inertia scale,
  \item SIM AHRS 기반 상태추정 안정화
\end{enumerate}

가 \textbf{겹쳐서} 만들어진 결과다.

\section{MuJoCo 공식 항과 현재 v3.0 코드의 대응관계}

\begin{center}
\small
\begin{tabularx}{\linewidth}{>{\raggedright\arraybackslash}p{0.18\linewidth}YY}
\toprule
공식 항 / 기능 & 현재 v3.0에서 담당 주체 & repository 기준 위치 \\
\midrule
ellipsoid added mass \(\mathbf{f}_A, \mathbf{g}_A\) & MuJoCo built-in ellipsoid fluid & \code{tank\_current\_scene.xml} 의 \code{fluidshape}, \code{fluidcoef} \\
quadratic drag \(\mathbf{f}_D\) & MuJoCo built-in ellipsoid fluid & same \\
angular drag \(\mathbf{g}_D\) & MuJoCo built-in ellipsoid fluid & same \\
linear viscous resistance \(\mathbf{f}_V, \mathbf{g}_V\) & MuJoCo built-in ellipsoid fluid + medium viscosity & XML \code{density}, \code{viscosity} \\
Magnus lift \(\mathbf{f}_M\) & 공식 모델엔 있음, current scene에선 \(C_M=0\) 이라 사실상 비활성 & \code{fluidcoef[4]=0.0} \\
Kutta lift \(\mathbf{f}_K\) & 공식 모델엔 있음, current scene에선 \(C_K=0\) 이라 사실상 비활성 & \code{fluidcoef[3]=0.0} \\
distributed buoyancy / CoB torque & Python runtime & \code{run\_urdf\_full.py}, \code{sim\_profiles.json} \\
body mass / inertia synthesis & Python runtime에서 합성 후 MuJoCo body inertial 반영 & \code{sim\_profiles.json} \code{body\_components} \\
thruster lag / deadzone / scaling & Python runtime & \code{run\_urdf\_full.py}, \code{thruster\_params.json} \\
\bottomrule
\end{tabularx}
\end{center}

\subsection{무엇을 MuJoCo에 맡기고 무엇을 runtime이 보정하는가}

이 부분을 잘못 설명하면 ``current가 built-in ellipsoid fluid model이다'' 라는 말이 과장된다. 정확한 설명은 아래와 같다.

\begin{itemize}[leftmargin=2em]
  \item \textbf{Fluid drag / added mass / lift 계산}: MuJoCo 공식 ellipsoid model
  \item \textbf{Hydrostatic buoyancy / restoring}: Python runtime
  \item \textbf{Thruster actuator shaping}: Python runtime
  \item \textbf{Rigid-body integration}: MuJoCo
  \item \textbf{Sensor truth export to ArduSub SITL}: Python runtime bridge
\end{itemize}

따라서 current path는 \textbf{MuJoCo fluid term을 주력으로 쓰되, vehicle-specific hydrostatic/actuation/interface는 runtime에서 유지하는 하이브리드 구조}다.

\section{Legacy와 current의 핵심 차이}

\subsection{Legacy는 무엇이었는가}

legacy path는 본질적으로 \textbf{Python runtime이 coefficient-based 6-DOF hydrodynamics를 직접 계산}하는 구조였다. 즉

\begin{itemize}[leftmargin=2em]
  \item added mass
  \item added-mass Coriolis
  \item linear damping
  \item quadratic damping
\end{itemize}

을 runtime helper가 body frame에서 계산해 \code{xfrc\_applied} 로 넣었다.

\subsection{Current는 무엇으로 바뀌었는가}

current path는 이 custom hydro force를 끄고, 대신

\begin{enumerate}[leftmargin=2em]
  \item scene에서 \code{fluidshape="ellipsoid"} 활성화
  \item geom별 \code{fluidcoef} 설정
  \item current profile로 drag / inertia / buoyancy / thruster authority 재구성
\end{enumerate}

구조로 바뀌었다.

\subsection{안정화에 가장 크게 기여한 것은 무엇인가}

\textbf{fluid force 관점만 놓고 보면} 가장 크게 작용한 변화는 다음 두 개다.

\begin{enumerate}[leftmargin=2em]
  \item \textbf{custom 6-DOF hydrodynamic wrench를 끄고 MuJoCo built-in ellipsoid fluid를 주력으로 쓴 것}
  \item \textbf{current profile에서 inertia scale을 극단적으로 낮추고 drag / angular drag baseline을 키운 것}
\end{enumerate}

첫 번째 변화는 force model을 더 shape-aware 하게 바꿨고, 두 번째 변화는 그 모델의 결과가 자세/속도에 더 빠르게 정리되도록 만들었다.

다만 \textbf{depth hold 안정화 전체}를 설명할 때는 fluid model만으로는 불충분하다. 실제 depth-hold 무한하강 문제를 끝낸 결정적 수정에는 \code{AHRS\_EKF\_TYPE=10}, \code{RNGFND1\_*} 정렬 같은 \textbf{상태추정/센서 backend 수정}이 포함된다. 즉 \textbf{fluid model 개선과 estimator 안정화가 함께 들어가야 현재 closed-loop가 성립한다.}

\section{각 항을 튜닝하면 어떤 현상이 바뀌는가}

\begin{center}
\small
\begin{tabularx}{\linewidth}{>{\raggedright\arraybackslash}p{0.22\linewidth}Y}
\toprule
조정 항목 & 먼저 나타나는 현상 \\
\midrule
\(\rho\) 증가 & 가속 저항, drag, lift 전반 증가. 물이 더 ``무거워짐'' \\
\(\beta\) 증가 & 저속 damping 증가, 잔류 속도와 잔진동 감소 \\
\(C_{D,\text{blunt}}\) 증가 & broadside drift, 측면 밀림, 수평 이동 저항 증가 \\
\(C_{D,\text{slender}}\) 증가 & 전진축/길이방향 저항 증가, angular drag 혼합 증가 \\
\(C_{D,\text{angular}}\) 증가 & yaw/roll/pitch tail motion 감소, 회전 정리 빨라짐 \\
\(C_K\) 증가 & 받음각이 있는 움직임에서 side-force / lift 커짐 \\
\(C_M\) 증가 & spin과 translation이 섞일 때 옆으로 휘는 coupling 커짐 \\
geom size 증가 & projected area, volume, added mass 모두 증가 \\
geom pos 이동 & 같은 힘이어도 CoM 기준 torque가 달라짐 \\
\bottomrule
\end{tabularx}
\end{center}

\section{이 모델의 한계}

공식 문서가 강조하듯 MuJoCo fluid model은 \textbf{state-less phenomenological model} 이다. 따라서 다음은 기대하면 안 된다.

\begin{itemize}[leftmargin=2em]
  \item wake shedding
  \item 자유수면 wave interaction
  \item pressure field 직접 해석
  \item 두 body 사이의 복잡한 유동 간섭
  \item 자유표면 근처의 세밀한 공기-물 혼상 효과
\end{itemize}

즉 current path가 legacy보다 더 안정적이고 설명력 있는 것은 맞지만, 이것이 CFD 대체물이라는 뜻은 아니다. 이 모델의 장점은 \textbf{실시간 제어기와 붙일 수 있을 만큼 빠르면서도, 방향성 drag/added-mass/lift를 보수적으로 표현할 수 있다는 점}이다.

\section{결론}

MuJoCo 공식 문서 기준으로 현재 v3.0 current path를 가장 정확하게 설명하면 다음 한 문장으로 정리된다.

\begin{quote}
현재 모델은 \textbf{MuJoCo built-in ellipsoid fluid model} 을 주력 fluid force 엔진으로 사용하고, 그 주변의 \textbf{buoyancy, distributed mass/inertia, thruster shaping, SITL sensor bridge} 를 Python runtime이 보정하는 \textbf{하이브리드 수중 시뮬레이션 구조}다.
\end{quote}

그리고 이 current path의 안정화는 단일 계수 하나에서 나온 것이 아니다. \textbf{공식 ellipsoid fluid force의 도입}, \textbf{current profile의 보수적 drag / inertia 설정}, \textbf{상태추정 경로의 SIM backend 정렬}이 동시에 맞물리면서 지금의 거동이 나온다.

\section*{공식 소스 링크}

\begin{itemize}[leftmargin=2em]
  \item MuJoCo official docs: \href{https://mujoco.readthedocs.io/en/3.1.1/computation/fluid.html}{Fluid forces}
  \item MuJoCo official docs: \href{https://mujoco.readthedocs.io/en/3.1.1/XMLreference.html}{XML reference}
\end{itemize}

\end{document}
